% Revize 2026-10-28: obsah, matematika a návaznost; podrobnosti a časy
% viz planovani/REVIZE_PRIPRAV_1_ROCNIK.md.
% Aktuální porovnání se skutečným obsahem místních PDF Realisticky.cz:
% externi_zdroje/realisticky/01_Zakladni_poznatky/08_Mnohocleny/010803_Pouziti_vzorcu_pri_uprave_mnohoclenu_I/010803_Lekce.pdf
% externi_zdroje/realisticky/01_Zakladni_poznatky/08_Mnohocleny/010804_Pouziti_vzorcu_pri_uprave_mnohoclenu_II/010804_Lekce.pdf
% externi_zdroje/realisticky/01_Zakladni_poznatky/08_Mnohocleny/010808_Rozklad_mnohoclenu_na_soucin_III/010808_Lekce.pdf
% Revize 2026-10-27: cílené doplnění; nové úlohy a řešení jsou vlastní tvorba.
% Metodický vzor: 4_rocnik/pripravy/Uvod_do_kombinatoriky.pdf (samostatný dokument).
% Hodina 51, 1. ročník SPŠ. Rozsah: Rozpis_vyuky_MAT_1_rocnik.tex, řádek 51.

% Plán: 1_rocnik/tematicke_plany/SPS.pdf; společná projektová šablona.

% Učebnice: ucebnice/01_Zakladni_poznatky/Algebraicke_vyrazy/Algebraické výrazy.tex

% Původ úloh: Procvičení c): učebnice, „Rozklad mnohočlenů“, příklad x^5+x^3-x^2-1. Rezerva: přesný výraz z kostry učitele. Ostatní vlastní.

% Metodický zdroj: externi_zdroje/realisticky/01_Zakladni_poznatky/08_Mnohocleny/010804_Pouziti_vzorcu_pri_uprave_mnohoclenu_II/010804_Lekce.pdf

% URL: http://www.realisticky.cz/ucebnice/01%20Matematika%20S%C5%A0/01%20Z%C3%A1kladn%C3%AD%20poznatky/08%20Mnoho%C4%8Dleny/04%20Pou%C5%BEit%C3%AD%20vzorc%C5%AF%20p%C5%99i%20%C3%BAprav%C4%9B%20mnoho%C4%8Dlen%C5%AF%20II.pdf

% Metodický zdroj: externi_zdroje/realisticky/01_Zakladni_poznatky/08_Mnohocleny/010806_Rozklad_mnohoclenu_na_soucin_I_vytykani/010806_Lekce.pdf

% URL: http://www.realisticky.cz/ucebnice/01%20Matematika%20S%C5%A0/01%20Z%C3%A1kladn%C3%AD%20poznatky/08%20Mnoho%C4%8Dleny/06%20Rozklad%20mnoho%C4%8Dlen%C5%AF%20na%20sou%C4%8Din%20I%20(vyt%C3%BDk%C3%A1n%C3%AD).pdf

\documentclass[11pt,a4paper]{article}

% Spolecny styl pro ucitelske pripravy MAT (Work/MAT)
% Verze 1.1, 2026-10-23
\NeedsTeXFormat{LaTeX2e}
\ProvidesPackage{priprava}[2026/10/23 v1.1 Sablona ucitelskych priprav MAT]

% Zakladni typografie a cestina
\RequirePackage[T1]{fontenc}
\RequirePackage[utf8]{inputenc}
\RequirePackage[czech]{babel}
\RequirePackage{lmodern}
\RequirePackage{microtype}
\RequirePackage{amsmath,amssymb}
\RequirePackage{xcolor}
\RequirePackage{enumitem}
\RequirePackage[most]{tcolorbox}
\RequirePackage{tabularx,booktabs,array}
\RequirePackage{fancyhdr}
\RequirePackage{needspace}
\RequirePackage{geometry}
\RequirePackage[hidelinks]{hyperref}

\geometry{a4paper,top=17mm,bottom=18mm,left=23mm,right=23mm}
\definecolor{notegray}{gray}{0.39}
\setlength{\parindent}{0pt}
\setlength{\parskip}{2.5pt plus 1pt minus 0.5pt}
\setlist[itemize]{topsep=2pt,itemsep=1pt,parsep=0pt,leftmargin=1.45em}
\setlist[enumerate]{topsep=2pt,itemsep=2pt,parsep=0pt,leftmargin=1.65em}
\setcounter{secnumdepth}{0}
\renewcommand{\arraystretch}{1.12}

% Zahlavi: vlevo rocnik | globalni tema, vpravo cislo: nazev hodiny
% Zapati zustava zcela prazdne.
\pagestyle{fancy}
\fancyhf{}
\renewcommand{\headrulewidth}{0.35pt}
\renewcommand{\footrulewidth}{0pt}
\setlength{\headheight}{14pt}
\newcommand{\nastavPripravu}[4]{%
  \fancyhead[L]{\small #1\enspace\textbar\enspace #2}%
  \fancyhead[R]{\small #3: #4}%
}

% Minimalisticke cernobile boxy. Kazda definice zvlast.
\tcbset{commonlesson/.style={enhanced,breakable=false,colback=white,colframe=black,
  colbacktitle=white,coltitle=black,boxrule=0.45pt,
  arc=0.8mm,outer arc=0.8mm,left=3mm,right=3mm,top=1.6mm,bottom=1.6mm,
  toptitle=1mm,bottomtitle=0.7mm,fonttitle=\bfseries,
  attach boxed title to top left={xshift=2.5mm,yshift=-1.15mm},
  boxed title style={colback=white,colframe=white,boxrule=0pt,left=1mm,right=1mm}}}
\newtcolorbox{zadaniBox}[1][]{commonlesson,title={Zadání},#1}
\newtcolorbox{prikladBox}[1][]{commonlesson,title={Příklad},#1}
\newtcolorbox{definiceBox}[1][]{commonlesson,title={Definice},#1}

% Spolecny box pro zadani a jeho reseni, bez mezery mezi castmi.
% Pouziti:
%   \begin{prikladReseniBox}
%     Zadani ulohy.
%     \tcblower
%     \begin{align*} ... \end{align*}
%   \end{prikladReseniBox}
% Spodni cast ma jemny sedy podklad a delici linku; lze v ni sazet
% vice radku matematiky. Obe casti maji jen jeden spolecny ramecek.
\tcbset{solutionlesson/.style={
  commonlesson,
  enhanced jigsaw,
  breakable,
  colbacklower=black!3,
  segmentation style={solid,draw=black!25,line width=0.35pt},
  fontlower=\small\color{notegray},
  before lower={\noindent\textbf{Řešení:}\enspace},
  leftlower=3mm,rightlower=3mm
}}
\newtcolorbox{prikladReseniBox}[1][]{solutionlesson,title={Příklad},#1}
\newtcolorbox{zadaniReseniBox}[1][]{solutionlesson,title={Zadání},#1}

% Pomucky pro autora pripravy: sedive poznamky; \reseni pro starsi pripravy
\newcommand{\poznamka}[1]{\par\begingroup\small\color{notegray}\itshape
  \textnormal{\textbf{Poznámka pro učitele.}}\enspace #1\par\endgroup}
\newcommand{\reseni}[1]{\par\begingroup\small\color{notegray}
  \textbf{Očekávané řešení:}\enspace #1\par\endgroup}
\newcommand{\blok}[1]{\needspace{5\baselineskip}\subsection*{#1}}
\newcommand{\vyrok}[1]{\uv{#1}}


\nastavPripravu{1. ročník}{Algebraické výrazy}{51}{Třetí mocniny a rozklady}

\begin{document}

\blok{Výklad a souvislosti}

Rozlišuj třetí mocninu celého dvojčlenu a součet třetích mocnin. Potřebné vztahy odvodíme násobením a potom je použijeme při rozkladu. Rozvoj součinu tří závorek navazuje na postup ze 44. hodiny.

\begin{definiceBox}[title={Definice},breakable=false]
\textbf{Třetí mocnina dvojčlenu} $(a+b)^3$ znamená součin tří stejných dvojčlenů. Výraz $a^3+b^3$ naproti tomu obsahuje jen dva sčítance.
\end{definiceBox}


Pro reálná $a,b$ platí
\begin{align*}
a^3-b^3&=(a-b)(a^2+ab+b^2),\\
a^3+b^3&=(a+b)(a^2-ab+b^2),\\
(a+b)^3&=a^3+3a^2b+3ab^2+b^3,\\
(a-b)^3&=a^3-3a^2b+3ab^2-b^3.
\end{align*}
První vztah ověř roznásobením a zrušením prostředních členů, druhý získej nahrazením $b$ číslem $-b$. Rozvoj krychle odvoď jako $(a+b)^2(a+b)$; u rozdílu opět nahraď $b$ číslem $-b$.

Například při ověření prvního vzorce skutečně vzniká
\[
(a-b)(a^2+ab+b^2)
=a^3+a^2b+ab^2-a^2b-ab^2-b^3=a^3-b^3.
\]
U třetí mocniny dvojčlenu se naopak smíšené členy nevyruší:
člen $a^2b$ vzniká třemi volbami toho, ze které závorky vezmeme $b$,
a stejně třikrát vznikne $ab^2$.
Proto se rozvoj $(a+b)^3$ liší od dvoučlenného součtu $a^3+b^3$.

\blok{Řešené příklady}

\begin{prikladReseniBox}[title={Příklad 1},breakable=false]
Rozlož $x^3-8$ a $8a^3+27b^3$.
\tcblower
$x^3-8=(x-2)(x^2+2x+4)$. Dále $8a^3+27b^3=(2a+3b)(4a^2-6ab+9b^2)$. Ve druhé závorce jsou smíšený součin a dvě druhé mocniny.
\end{prikladReseniBox}


\begin{prikladReseniBox}[title={Příklad 2},breakable=false]
Roznásob $(2x-1)^3$.
\tcblower
$8x^3-12x^2+6x-1$. Čtyři členy vzniknou jako $(2x)^3-3(2x)^2\cdot1+3(2x)\cdot1^2-1^3$.
\end{prikladReseniBox}


\blok{Samostatné procvičení}

\begin{zadaniReseniBox}[title={Procvičení},breakable=false]
\begin{enumerate}[label=\alph*)]
\item Rozlož $27x^3-1$.
\item Roznásob $(x+2)^3$.
\item Rozlož $x^5+x^3-x^2-1$.
\end{enumerate}
\tcblower
\begin{enumerate}[label=\alph*)]
\item $(3x-1)(9x^2+3x+1)$.
\item $x^3+6x^2+12x+8$.
\item $x^3(x^2+1)-(x^2+1)=(x^2+1)(x^3-1)=(x^2+1)(x-1)(x^2+x+1)$.
\end{enumerate}
\end{zadaniReseniBox}


\poznamka{Hodina propojuje čtyři vzorce a je obsahově náročná. Základ tvoří jejich odvození, P1 a procvičení a), b); P2 lze rozebrat krátce jako připomenutí násobení. Procvičení c) a náročnou rezervu ponech rychlejším žákům nebo domů. Nepřidávej další hodinu ani nové vzorce.}

\blok{Volitelná rezerva}

\begin{prikladReseniBox}[title={Rozšíření procvičení},breakable=false]
Rozlož výraz ze sešitu $3x^6-2x^3-3x^4+2$ na součin.
\tcblower
Seskupíme $3x^4(x^2-1)-2(x^3-1)$ a použijeme oba vzorce:
\[
(x-1)\bigl[3x^4(x+1)-2(x^2+x+1)\bigr].
\]
Po úpravě hranaté závorky dostaneme $3x^5+3x^4-2x^2-2x-2$. Dělení tohoto mnohočlenu výrazem $x-1$ má nulový zbytek a dává $3x^4+6x^3+6x^2+4x+2$. Výsledný součin je
\[
(x-1)^2(3x^4+6x^3+6x^2+4x+2).
\]
Správnost ověříme roznásobením; není třeba tvrdit, že poslední činitel už nelze dále rozložit.
\end{prikladReseniBox}



\poznamka{Původní obtížnost i náročný příklad ze sešitu zachováme. Základ a rezervy jsou rozdělené v původní poznámce; nová domácí a) až c) slouží upevnění, d) je výslovně dobrovolná rezerva. Při rychlém zvládnutí odvození přidej původní procvičení c), nikoli další vzorce. Všechny čtyři vzorce za jednu hodinu jsou náročné i při svižném tempu; nevyžaduj současně všechny rezervy.}

\blok{Domácí procvičení}

\begin{zadaniReseniBox}[title={Domácí úlohy: základ a jedna rezerva},breakable=false]
\begin{enumerate}[label=\alph*)]
\item Rozlož $x^3+27$ a $8a^3-1$.
\item Roznásob $(x-2)^3$.
\item Rozlož $2x^4-16x$.
\item Rezerva: rozlož $x^6-1$ s využitím rozdílu čtverců i třetích mocnin.
\end{enumerate}
\tcblower
\begin{enumerate}[label=\alph*)]
\item $(x+3)(x^2-3x+9)$ a $(2a-1)(4a^2+2a+1)$.
\item $x^3-6x^2+12x-8$.
\item $2x(x^3-8)=2x(x-2)(x^2+2x+4)$.
\item $(x^3-1)(x^3+1)=(x-1)(x^2+x+1)(x+1)(x^2-x+1)$.
\end{enumerate}
\end{zadaniReseniBox}

\end{document}
